\chapter{Explorer}
\label{ch:explorer}
% BUDGET : 7 pages (cf. memoire/PLAN.md, section 4).
% THESE DU CHAPITRE : un classement de politiques publiques n'a de sens que relativement a un
% forcage, et le dispositif qui le produit doit enoncer ce qu'il ne peut PAS trancher.
%
% NOTE : les trois tableaux (grille, front budget, realloc) sont CONSERVES -- ils sont
% entierement bases sur des macros generees, il n'y a rien a retaper.

\minisommaire{%
  \ms{sec:dispositif}{politiques, forçages, balayages, et les trois décisions de méthode}%
  \ms{sec:grille}{le classement des politiques dépend du forçage}%
  \ms{sec:fronts}{ce que coûtent un euro public et un kilogramme d'azote}}

Le modèle est porté et calibré. On s'intéresse désormais à l'exploration de scénarios. Ce
chapitre présente le dispositif et les résultats.

\section{Le dispositif expérimental}
\label{sec:dispositif}

Chaque scénario est séparé en trois rubriques :

\begin{itemize}[leftmargin=1.2em,itemsep=6pt]
  \item La politique, qui correspond à une décision de la puissance publique. Cela peut
    être un plafond d'azote, une enveloppe de subventions, un plancher d'emploi, etc.
  \item Le forçage, qui correspond à ce que le climat et les marchés imposent. Cela peut
    être des prix, des rendements, des coûts d'intrants multipliés par un facteur, etc.
  \item Le balayage, qui correspond à une variation d'un seul seuil afin de tracer le
    front des compromis atteignables.
\end{itemize}

Le catalogue actuel est composé de \num{\mesPolitiques} politiques, \num{\mesForcages} forçages
et \num{\mesBalayages} fronts.

La figure~\ref{fig:spectre} donne les deux catalogues. Le catalogue des politiques est ordonné
d'une dérégulation complète à une rupture agroécologique. On peut voir quels scénarios ont été
explorés. Le produit cartésien complet donne \num{\mesCellulesTotal} combinaisons ; à la durée
moyenne observée en configuration prospective (\SI{\mesSolveProsMoy}{\second}, §~\ref{sec:tractabilite}),
cela représente une dizaine de jours de calcul continu, sans compter les fronts. Un plan
explicite désigne donc, pour chaque politique, ses forçages et ses fronts.

\begin{figure}[h]
  \centering
  % Une ligne de catalogue = un rectangle plein (la teinte porte le rang), trois noeuds :
  % marqueur de calcul, abreviation, nom.
  %   #1 x du panneau  #2 y du haut de la ligne  #3 teinte  #4 marqueur  #5 abrev.  #6 nom
  \newcommand{\scenrow}[6]{%
    \fill[#3] (#1,#2) rectangle ++(69,-4.2);
    \node[anchor=west, inner sep=0] at (#1+2.2,#2-2.1) {#4};
    \node[anchor=west, inner sep=0, font=\scriptsize\bfseries] at (#1+5.8,#2-2.1) {#5};
    \node[anchor=west, inner sep=0] at (#1+14,#2-2.1) {#6};}
  \newcommand{\calc}{$\bullet$}%   au moins une cellule resolue
  \newcommand{\jamais}{$\circ$}%   jamais calcule
  \begin{tikzpicture}[
      x=1mm, y=1mm, font=\scriptsize,
      % Les quatre encadres se distinguent par la COULEUR *et* par le trait : imprimee en
      % noir et blanc, la figure reste lisible.
      temoin/.style  = {draw=black,       line width=0.9pt, rounded corners=0.7mm},
      climat/.style  = {draw=cadreClimat, line width=0.8pt, rounded corners=0.7mm},
      marche/.style  = {draw=cadreMarche, line width=0.8pt, rounded corners=0.7mm,
                        dash pattern=on 1.4mm off 0.9mm},
      horsgrad/.style= {draw=black!65,    line width=0.8pt, rounded corners=0.7mm,
                        dash pattern=on 0.35mm off 0.7mm}]

    % ================= PANNEAU GAUCHE : LES DOUZE FORÇAGES =================
    % Ordre : les deux bornes en tete, les deux blocs thematiques au milieu (c'est eux
    % qu'on encadre, donc leurs lignes doivent etre contigues), la borne basse en pied.
    % A l'interieur d'un bloc, les lignes vont du moins au plus severe.
    \node[anchor=south west, font=\scriptsize] at (1,-0.3)
      {Forçages --- ce qu'imposent le climat et les marchés};
    \draw[black!45, line width=0.3pt] (1,-0.9) -- (70,-0.9);

    \scenrow{1}{-1.5} {scenVert!75!scenSable}{\jamais}{F10}{conjoncture favorable}
    \scenrow{1}{-6.1} {scenSable}            {\calc}  {F0} {nominal (forçage nul)}

    \scenrow{1}{-12.5}{scenRouge!30!scenSable}{\jamais}{F3}{dérive climatique}
    \scenrow{1}{-17.1}{scenRouge!35!scenSable}{\jamais}{F4}{choc sanitaire}
    \scenrow{1}{-21.7}{scenRouge!55!scenSable}{\calc}  {F1}{cyclone majeur}
    \scenrow{1}{-26.3}{scenRouge!60!scenSable}{\jamais}{F2}{sécheresse sévère}

    \scenrow{1}{-32.7}{scenVert!35!scenSable} {\calc}  {F7}{inflation alimentaire}
    \scenrow{1}{-37.3}{scenRouge!40!scenSable}{\jamais}{F8}{retrait du \glsentryshort{posei}}
    \scenrow{1}{-41.9}{scenRouge!45!scenSable}{\jamais}{F5}{effondrement des prix d'export}
    \scenrow{1}{-46.5}{scenRouge!50!scenSable}{\calc}  {F6}{choc sur les intrants}

    \scenrow{1}{-53.1}{scenRouge!30!scenSable}{\jamais}{F11}{artificialisation des terres}
    \scenrow{1}{-57.7}{scenRouge!80!scenSable}{\calc}  {F9} {crise systémique}

    \draw[temoin] (0.1,-5.3)  rectangle (70.9,-11.1);
    \draw[climat] (0.1,-11.7) rectangle (70.9,-31.3);
    \draw[marche] (0.1,-31.9) rectangle (70.9,-51.5);

    % ================= PANNEAU DROIT : LES DIX POLITIQUES =================
    \node[anchor=south west, font=\scriptsize] at (80,-0.3)
      {Politiques --- ce que décide la puissance publique};
    \draw[black!45, line width=0.3pt] (80,-0.9) -- (149,-0.9);

    \scenrow{80}{-1.5} {scenVert!75!scenSable} {\jamais}{P10}{bifurcation agroécologique}
    \scenrow{80}{-6.1} {scenVert!58!scenSable} {\calc}  {P8} {transition agroécologique}
    \scenrow{80}{-10.7}{scenVert!42!scenSable} {\calc}  {P7} {Écophyto réglementaire}
    \scenrow{80}{-15.3}{scenVert!25!scenSable} {\calc}  {P6} {verdissement incitatif}
    \scenrow{80}{-19.9}{scenSable}             {\calc}  {P4} {statu quo}
    \scenrow{80}{-26.3}{scenRouge!25!scenSable}{\jamais}{P3} {intensification modérée}
    \scenrow{80}{-30.9}{scenRouge!42!scenSable}{\calc}  {P5} {austérité budgétaire}
    \scenrow{80}{-35.5}{scenRouge!58!scenSable}{\jamais}{P2} {accélérationnisme industriel}
    \scenrow{80}{-40.1}{scenRouge!75!scenSable}{\calc}  {P1} {dérégulation totale}
    \scenrow{80}{-47.5}{scenGris}              {\calc}  {P9} {souveraineté alimentaire}

    \draw[temoin]   (79.1,-19.1) rectangle (149.9,-24.9);
    \draw[horsgrad] (79.1,-46.7) rectangle (149.9,-52.5);

    % ================= LÉGENDE =================
    % Le degrade se lit comme une bande : sept pastilles, du vert au rouge.
    \fill[scenVert!75!scenSable]  (1,-65.5)  rectangle ++(5,-2.9);
    \fill[scenVert!50!scenSable]  (6,-65.5)  rectangle ++(5,-2.9);
    \fill[scenVert!25!scenSable]  (11,-65.5) rectangle ++(5,-2.9);
    \fill[scenSable]              (16,-65.5) rectangle ++(5,-2.9);
    \fill[scenRouge!25!scenSable] (21,-65.5) rectangle ++(5,-2.9);
    \fill[scenRouge!50!scenSable] (26,-65.5) rectangle ++(5,-2.9);
    \fill[scenRouge!75!scenSable] (31,-65.5) rectangle ++(5,-2.9);
    \node[anchor=west, inner sep=0] at (38,-67)
      {teinte : du plus favorable, ou du plus agroécologique, au plus défavorable};

    \draw[temoin]   (1,-71.5)   rectangle ++(6,-2.9);
    \node[anchor=west, inner sep=0] at (8.5,-73) {témoin (F0, P4)};
    \draw[climat]   (32,-71.5)  rectangle ++(6,-2.9);
    \node[anchor=west, inner sep=0] at (39.5,-73) {climat et sanitaire};
    \draw[marche]   (69,-71.5)  rectangle ++(6,-2.9);
    \node[anchor=west, inner sep=0] at (76.5,-73) {marchés et politique};
    \draw[horsgrad] (107,-71.5) rectangle ++(6,-2.9);
    \node[anchor=west, inner sep=0] at (114.5,-73) {hors gradient};

    \node[anchor=west, inner sep=0] at (1,-78)
      {\calc~au moins une cellule de la grille a été résolue \qquad
       \jamais~aucune : la ligne n'a jamais été calculée};
  \end{tikzpicture}
  \caption{Les deux catalogues. La teinte va de la
  conjoncture la plus favorable (F10) à la plus défavorable (F9), et de la rupture
  agroécologique (P10) à la dérégulation (P1). Les témoins encadrés en noir
  sont les points de comparaison. F0 est le forçage nul, P4 le statu quo.}
  \label{fig:spectre}
\end{figure}

Trois décisions de méthode gouvernent ce dispositif.

\begin{enumerate}[leftmargin=1.6em,itemsep=8pt]

\item Les seuils du catalogue sont des fractions choisies de totaux mesurés sur un run de
référence, et non des chiffres inventés. L'annexe~\ref{ann:scenarios} donne, pour chaque seuil,
la fraction et le total dont elle est tirée.

\item On décrit une politique par ses résultats sous chaque forçage. On mobilise les 
critères classiques de décision dans l'incertain : pire cas \parencite{wald1945}, 
regret maximal \parencite{savage1951}, domaine de viabilité \parencite{starr1963}. L'infaisabilité correspond au pire résultat.

\item Les fronts sont tracés par $\varepsilon$-contrainte \parencite{haimes1971,mavrotas2009} et
non par somme pondérée. La raison est qu'une somme pondérée ne peut atteindre que les points
situés sur l'enveloppe convexe du front.

\end{enumerate}



\section{Ce que la grille montre}
\label{sec:grille}

Le classement des politiques dépend du forçage. Le tableau~\ref{tab:grille} donne les cellules
résolues ; l'annexe~\ref{ann:resultats} en donne les autres indicateurs, les durées et les
drapeaux de convergence.

Sous crise systémique (F9), les trois politiques testées donnent une marge brute du territoire
négative. C'est le pire cas au sens de Wald.

La politique de souveraineté alimentaire (P9) est la meilleure sous le forçage nominal (F0).
Elle donne \SI{\grillePneufFzero}{\mega\euros} de marge brute contre
\SI{\grillePquatreFzero}{\mega\euros} pour le statu quo (P4). Elle est également la pire sous
crise systémique (F9), avec \SI{\grillePneufFneuf}{\mega\euros} contre
\SI{\grillePquatreFneuf}{\mega\euros}. Le classement s'inverse d'un forçage à l'autre. Le regret
de Savage montre cette inversion (figure~\ref{fig:regret}). Passer de F0 à F9 fait passer le
regret de P9 de \SI{\regretPneufFzero}{\mega\euros} à \SI{\regretPneufFneuf}{\mega\euros}. À
l'inverse, le regret de P4 passe de \SI{\regretPquatreFzero}{\mega\euros} à
% CORRIGE : la phrase disait « de 30 % a 10 % ». Sous F9, P4 est la MEILLEURE politique, son
% regret y vaut donc 0. Les 9,82 M EUR sont son regret sous F6, pas sous F9.
% CORRIGE 2026-09-02 : l'unite etait le POURCENT. build_chiffres.py calcule ces macros comme
% (meilleure - valeur) / 1e6, donc en MILLIONS D'EUROS. Sous F9 la meilleure marge est
% negative (-16,99 M EUR), et un regret rapporte a une valeur negative n'a aucun sens.
\SI{\regretPquatreFneuf}{\mega\euros}.

Un point important est que la marge brute contient la subvention. On cherche à classer des
politiques et donc à classer notamment des déplacements d'aides. Regarder la marge brute, c'est,
en partie, avantager les politiques qui subventionnent davantage. On propose donc une colonne
\enquote{hors aide} qui isole la marge issue de la production agricole seule. On constate que le
classement change. Sous F9, P8 (\SI{\grilleNetPhuitFneuf}{\mega\euros}) est plus exposée que P9
(\SI{\grilleNetPneufFneuf}{\mega\euros}). Les métriques \enquote{brute} et \enquote{hors aide}
se complètent et ne donnent pas les mêmes informations.

Ce qui limite les politiques, c'est la contradiction entre hausse des \gls{etp} et régulation
des pratiques. Les itinéraires les plus intensifs en travail sont aussi les plus intensifs en
azote et en \gls{ift}. En somme, créer du travail peut devenir impossible à cause des autres
plafonds.

À noter qu'un forçage peut rendre une politique inutile, c'est-à-dire qu'elle ne coûte et
n'apporte rien. Les contraintes imposées par la politique peuvent être déjà contenues dans les
contraintes imposées par les forçages. Sous crise systémique, le plafond d'azote de P8 ne mord
jamais car la crise réduit les surfaces intensives.

\begin{table}[h]
  \centering
  \caption{Marge brute par cellule (M\eur).
  hors aide exclut les subventions de la marge brute. La dague marque un solve sous-optimal car arrêté à la limite de temps.}
  \label{tab:grille}
  \small
  \begin{tabularx}{\textwidth}{@{}>{\raggedright\arraybackslash}Xrrrrrr@{}}
    \toprule
    & \multicolumn{2}{c}{\textbf{F0 nominal}} & \multicolumn{2}{c}{\textbf{F6 intrants}}
    & \multicolumn{2}{c}{\textbf{F9 crise}} \\
    \cmidrule(lr){2-3}\cmidrule(lr){4-5}\cmidrule(lr){6-7}
    & brute & hors aide & brute & hors aide & brute & hors aide \\
    \midrule
    P4 statu quo & \num{\grillePquatreFzero}\grillePquatreFzeroFlag & \num{\grilleNetPquatreFzero}
                 & \num{\grillePquatreFsix}\grillePquatreFsixFlag   & \num{\grilleNetPquatreFsix}
                 & \num{\grillePquatreFneuf}\grillePquatreFneufFlag & \num{\grilleNetPquatreFneuf} \\
    P6 verdissement & \num{\grillePsixFzero}\grillePsixFzeroFlag & \num{\grilleNetPsixFzero}
                 & \num{\grillePsixFsix}\grillePsixFsixFlag & \num{\grilleNetPsixFsix}
                 & --- & --- \\
    P7 Écophyto  & \num{\grillePseptFzero}\grillePseptFzeroFlag & \num{\grilleNetPseptFzero}
                 & \num{\grillePseptFsix}\grillePseptFsixFlag & \num{\grilleNetPseptFsix}
                 & --- & --- \\
    P8 agroécologie & \num{\grillePhuitFzero}\grillePhuitFzeroFlag & \num{\grilleNetPhuitFzero}
                 & --- & ---
                 & \num{\grillePhuitFneuf}\grillePhuitFneufFlag & \num{\grilleNetPhuitFneuf} \\
    P9 souveraineté & \num{\grillePneufFzero}\grillePneufFzeroFlag & \num{\grilleNetPneufFzero}
                 & --- & ---
                 & \num{\grillePneufFneuf}\grillePneufFneufFlag & \num{\grilleNetPneufFneuf} \\
    \bottomrule
  \end{tabularx}
\end{table}

\begin{figure}[h]
  \centering
  \begin{tikzpicture}[x=1.25mm, y=1mm, font=\scriptsize]
    % Regret en % de la meilleure politique du MEME forcage. Barres empilees par politique :
    % F0 en fonce, F6 en gris moyen, F9 en gris clair. Un regret nul est marque par un trait.
    \foreach \g in {0,10,20,30,40,50} {
      \draw[black!12] (\g,2) -- (\g,-46);
      \node[below, inner sep=1.5pt] at (\g,-46) {\g};
    }
    \draw[black!50] (0,2) -- (0,-46);
    % --- P4 statu quo
    \node[left, inner xsep=3pt] at (0,-2) {P4 statu quo};
    \fill[black!70] (0,0)    rectangle (\regretPquatreFzero,-1.6);
    \fill[black!40] (0,-2)   rectangle (\regretPquatreFsix,-3.6);
    \fill[black!15] (0,-4)   rectangle (0.2,-5.6);
    % --- P6 verdissement
    \node[left, inner xsep=3pt] at (0,-11) {P6 verdissement};
    \fill[black!70] (0,-9)   rectangle (\regretPsixFzero,-10.6);
    \fill[black!40] (0,-11)  rectangle (0.2,-12.6);
    % --- P7 Ecophyto
    \node[left, inner xsep=3pt] at (0,-20) {P7 Écophyto};
    \fill[black!70] (0,-18)  rectangle (\regretPseptFzero,-19.6);
    \fill[black!40] (0,-20)  rectangle (\regretPseptFsix,-21.6);
    % --- P8 agroecologie
    \node[left, inner xsep=3pt] at (0,-29) {P8 agroécologie};
    \fill[black!70] (0,-27)  rectangle (\regretPhuitFzero,-28.6);
    \fill[black!15] (0,-29)  rectangle (\regretPhuitFneuf,-30.6);
    % --- P9 souverainete
    \node[left, inner xsep=3pt] at (0,-38) {P9 souveraineté};
    \fill[black!70] (0,-36)  rectangle (0.2,-37.6);
    \fill[black!15] (0,-38)  rectangle (\regretPneufFneuf,-39.6);
    % --- Legende
    \fill[black!70] (0,-52)  rectangle (2.5,-53.6); \node[right] at (2.5,-52.8) {F0 nominal};
    \fill[black!40] (16,-52) rectangle (18.5,-53.6); \node[right] at (18.5,-52.8) {F6 intrants};
    \fill[black!15] (33,-52) rectangle (35.5,-53.6); \node[right] at (35.5,-52.8) {F9 crise};
    \node[right] at (48,-52.8) {M\eur};
  \end{tikzpicture}
  \caption{Regret par politique et par forçage, en millions d'euros de marge brute perdus par
  rapport à la meilleure politique du même forçage. Une barre nulle correspond au meilleur choix
  de politique, sachant le forçage. Aucune politique n'a une barre nulle sous tous les forçages.}
  \label{fig:regret}
\end{figure}

\section{Le coût marginal : euro public et azote}
\label{sec:fronts}

\subsection*{L'enveloppe publique}

La marge brute varie en sens inverse du plafond de subventions sur les sept points
(tableau~\ref{tab:front-budget}). Par exemple, si on réduit l'enveloppe publique de
\SI{\frontsubvQuatreSeuilHaut}{\euros} à \SI{\frontsubvQuatreunSeuil}{\euros}, la marge brute du
territoire augmente de plus de \SI{14}{\mega\euros} pendant que l'objectif diminue de
\SI{\frontsubvQuatreObjHaut}{\mega\euros} à \SI{\frontsubvQuatreObjBas}{\mega\euros}. On rappelle
que l'objectif est une marge ajustée du risque. Cet ajustement correspond à \SI{18}{\percent} de
la marge brute pour le plafond haut de subventions et \SI{34}{\percent} pour le plafond bas de
subventions. On constate que le modèle va chercher des cultures plus rémunératrices et exposées.
L'aide publique achète la stabilité et non la marge.

On considère maintenant les points de contrôle du balayage de l'enveloppe publique afin de
valider le front. \num{\frontsubvQuatrePlateau} des \num{\frontsubvQuatrePoints} points ont un
plafond supérieur à ce que le statu quo dépense spontanément, donc ils ne contraignent rien.
L'objectif retombe exactement sur le run non balayé (\SI{\frontsubvQuatreTemoin}{\mega\euros}).

Le prix dual n'est a priori valable qu'au voisinage du point. L'intérêt d'un front est de
pouvoir visualiser l'évolution de ce prix sur un intervalle.

La différence entre le statu quo avec un plafond de subvention libre et le plafond le plus serré
est claire. La banane d'exportation passe de \SI{2004}{\ha} à \SI{95}{\ha}
(\SI{-95}{\percent}), la canne diminue de \SI{2231}{\ha}, le maraîchage et l'igname doublent, les
agrumes passent de \SI{9}{\ha} à \SI{657}{\ha}, la prairie gagne \SI{1539}{\ha}
(tableau~\ref{tab:ann-realloc}). Retirer l'aide ne transforme donc pas l'export en vivrier
hectare pour hectare : l'export libère \SI{4140}{\ha} (\SI{1909}{\ha} de banane, \SI{2231}{\ha}
de canne), et le vivrier n'en reprend qu'environ la moitié.

Le modèle dit que l'enveloppe publique maintient la banane d'exportation sur le sol
guadeloupéen et joue un rôle important dans l'existence de la culture de canne. En revanche, il
ne vérifie pas que les \SI{1092}{\ha} de maraîchage supplémentaires trouvent un acheteur
(chapitre~\ref{ch:discussion}).

\begin{table}[h]
  \centering
  \caption{Front marge $\times$ enveloppe publique sous P4, en euros. L'astérisque marque les
  points de contrôle : l'objectif y retombe exactement sur le run non balayé
  (\num{\frontsubvQuatreTemoin}~M\eur), donc le plafond n'y contraint rien.}
  \label{tab:front-budget}
  \small
  % Le tabular entier est genere (entete et filets compris) : le \input de LaTeX ouvre une
  % cellule et ne peut donc pas vivre A L'INTERIEUR d'un alignement.
  % Genere par memoire/build_chiffres.py -- ne pas editer.
\begin{tabular}{@{}lrr@{}}
  \toprule
  \textbf{Plafond} & \textbf{Niveau atteint} & \textbf{Objectif (M\eur)} \\
  \midrule
  \num{35900000} & \num{35898271} & \num{74.36} \\
  \num{43000000} & \num{42986537} & \num{76.84} \\
  \num{50300000} & \num{50070701} & \num{77.96} \\
  \num{57400000} & \num{57111266} & \num{79.18} \\
  \num{64600000} & \num{64590701} & \num{80.62} \\
  \num{72000000} & \num{71277791} & \num{81.04}\textsuperscript{*} \\
  \num{90000000} & \num{71277791} & \num{81.04}\textsuperscript{*} \\
  \bottomrule
\end{tabular}

\end{table}

\subsection*{L'azote}

Le second front que l'on considère est celui de l'azote. Il est tracé sous la configuration de
% CORRIGE : « sous le statu quo (P4) » -- le front de reference est trace sous la CALIBRATION
% RETENUE, dont l'objectif temoin vaut 81,22 M EUR, et non sous P4, dont l'objectif vaut 81,04.
référence, c'est-à-dire la calibration retenue, et sous la politique agroécologique (P8). La
comparaison des deux fronts donne comme résultat une pente moyenne de
\SI{\frontazoteRefPente}{\euros\per\kilogram} sous la référence contre
\SI{\frontazoteHuitPente}{\euros\per\kilogram} sous P8. Cela veut dire qu'un kilogramme d'azote
est moins cher à retirer sous une politique déjà contraignante. Ici, P8 impose un plafond
d'\gls{ift}, un plafond d'eau, un plancher d'emploi et une inertie de \SI{40}{\percent} par
rapport à l'assolement existant. Cela veut dire que les usages les plus intensifs sont déjà
écartés et que les hectares faciles à désazoter ont déjà été pris.

Sans tracer les fronts, on serait passé à côté de ce résultat car les duaux au point le plus
serré des deux configurations sont de \SI{\frontazoteRefDualBas}{\euros\per\kilogram} et
\SI{\frontazoteHuitDualBas}{\euros\per\kilogram}. On aurait conclu à l'égalité des coûts
marginaux, ce qui est réfuté par les fronts.
